% ─────────────────────────────────────────────────────────────────────────────
% NovelGym: Additional Implementation Details
% Included inline at the end of the NovelGym section in Chapter 3
% ─────────────────────────────────────────────────────────────────────────────

\subsection{Convergence Criteria}
\label{sec:ng_convergence}

\rev{The learner is considered to have converged if the following conditions are met:}
\begin{itemize}
    \item \rev{The agent achieves the goal \( \delta g \geq \delta G \) in the last \( \eta \) epochs.}
    \item \rev{The average reward in these epochs is \( \delta r \geq \delta R \).}
    \item \rev{The max success rate and the max reward in the past \( \eta + \upsilon \) epochs is not greater to the those in the last \( \eta \) epochs.}
\end{itemize}
\rev{where each epoch is 4800 time steps, $\eta = 5$, $\upsilon = 5$, $\delta G = 0.9$, $\delta R = 400$. The episode length is 400 time steps. Episode ends if the agent reaches the goal state or if it completes 400 time steps.}


\subsection{Action Space}
\label{sec:ng_action_space}

\rev{The action space for the hybrid agent is composed of higher-level action operators such as \texttt{approach <entity>}. These are implemented as executors that use the primitive actions available in the environment to execute the operators. Therefore the hybrid agent's action space includes both higher-level operators (implemented as lower-level action executors) and primitive actions. For the RL agent, each higher-level operator is grounded. The exhaustive list of the actions used:}
\begin{itemize}
    \item \rev{\texttt{collect}}
    \item \rev{\texttt{break\_block}}
    \item \rev{\texttt{approach\_oak\_log}}
    \item \rev{\texttt{approach\_diamond\_ore}}
    \item \rev{\texttt{approach\_crafting\_table}}
    \item \rev{\texttt{approach\_block\_of\_platinum}}
    \item \rev{\texttt{approach\_entity\_103}}
    \item \rev{\texttt{interact\_103}}
    \item \rev{\texttt{select\_oak\_log}}
    \item \rev{\texttt{select\_iron\_pickaxe}}
    \item \rev{\texttt{select\_sapling}}
    \item \rev{\texttt{select\_tree\_tap}}
    \item \rev{\texttt{select\_crafting\_table}}
    \item \rev{\texttt{deselect\_item}}
    \item \rev{\texttt{craft\_stick}}
    \item \rev{\texttt{craft\_planks}}
    \item \rev{\texttt{craft\_block\_of\_diamond}}
    \item \rev{\texttt{craft\_pogo\_stick}}
    \item \rev{\texttt{trade\_block\_of\_titanium\_1}}
    \item \rev{\texttt{move\_forward}, \texttt{move\_backward}, \texttt{move\_left}, \texttt{move\_right}}
    \item \rev{\texttt{rotate\_left}, \texttt{rotate\_right}}
    \item \rev{\texttt{place}}
    \item \rev{\texttt{<slot 1>}, \texttt{<slot 2>} (vary during post-novelty situations)}
\end{itemize}


\subsection{Observation Spaces}
\label{sec:ng_observation_spaces}

\subsubsection{LiDAR Representation}

\rev{In this representation, the agent's observation space is a Gymnasium Box space. This Box space is made from a one-dimensional vector created by concatenating three vectors for the world as observed by the agent, the agent's inventory, and the item selected by the agent. The agent's world observation is generated in the following way: 8 beams of $45^\circ$ each are sent from $0^\circ$ to $360^\circ$, and the Euclidean distances of the objects that strike the LiDAR are stored in a two-dimensional array whose rows correspond to the integer encodings of all the possible objects and entities in the environment and whose columns correspond to the individual directions in which a LiDAR beam is sent. The assumption is made that occlusions do not hold true. The two-dimensional array is subsequently flattened. The inventory is represented as a one-dimensional array where the integer at each index corresponds to the number of objects encoded by the index that the agent has. The selected item is represented as a one-hot encoding in a one-dimensional array. Objects and entities in all three cases are assigned integer encodings. The low and high bounds of the Box space are arrays of a length corresponding to the maximum number of object and entity types in the environment and of maximum values corresponding to the maximum count for each object and entity type. The size of the observation used during our experiment is 250, which includes placeholders for two additional objects in post-novelty situations. We added placeholders for convenience in implementation. In principle our architecture can automatically expand the network based on new entities.}

\subsubsection{Image-based Representation}

\rev{In this representation, the observation space of the agent is a dictionary with three key-value pairs: for the agent's local view, for the agent's inventory, and for the agent's selected item. Each of the values in the dictionary is a Gymnasium Box space, where the low and high bounds are given by the maximum number of instances of an object or entity in the environment. The agent's local view is represented as a two-dimensional square image, where the number of one-hot encoded channels corresponds to the number of possible object and entity types in the environment. The inventory is represented as a one-dimensional array where the integer at each index corresponds to the number of objects encoded by the index that the agent has. The selected item is represented as a one-hot encoding in a one-dimensional array. Objects and entities in all three cases are represented by integer encodings.}


\subsection{Novelty Implementations}
\label{sec:ng_novelty_implementations}

\rev{For a list of the environment transformations currently implemented in the NovelGym environment, see Table~\ref{tab:rl_systemic_examples}. The components of a novelty being integrated into the system are as follows.}
\begin{itemize}
    \item \rev{\textbf{YAML config file}: This config file serves as an add-on to the existing YAML config file that configures the environment. Anything included in the new config expands on the settings in the environment config, and anything included under the key ``novelties'' overrides any chosen component of the original settings. Common entries to the config add or modify actions, action sets, action preconditions and effects, entities, entity properties, and the grid world layout. A novelty can be created through a YAML config file alone when it does not involve adding or modifying any modules, an example of which are the \textit{Multi-room} and \textit{Distance} novelties in Table~\ref{tab:rl_systemic_examples}. This is possible when no new source Python module is required for the novelty to function.}
    \item \rev{\textbf{Python source file(s)}: In the case that a novelty requires the addition or modification of an entity or action module, it is necessary to include the Python source code for this adjustment. All the novelties specified in Table~\ref{tab:rl_systemic_examples} except for \textit{Multi-room} and \textit{Distance} involve the addition of one or more modules. Note that changing the details of an existing module can be accomplished by changing the source Python code for the module in the novelty config file.}
\end{itemize}

\begin{table}[!tb]
    \centering
    \footnotesize
    \begin{tabular}{p{2.2cm} p{2.2cm} p{9cm}}
    \toprule
    
    \textbf{Novelty} & \textbf{Category}  & \textbf{Description} \tabularnewline  \arrayrulecolor{taupegray} \midrule
    
    \emph{Axe} & \emph{Detrimental} & The agent is provided with an axe in its inventory at the start of the episode, and a new action ``select axe'' is added to its action set. The agent cannot break a tree unless it has selected the axe beforehand. \tabularnewline
    \arrayrulecolor{taupegray} \midrule
    
    \emph{Busy} & \emph{Nuisance} & Each trader is busy during a certain percentage of time steps. During time steps that a trader is busy, the agent cannot successfully execute a trade against them. \tabularnewline
    \arrayrulecolor{taupegray} \midrule
    
    \emph{Chest} & \emph{Beneficial} & A plastic chest containing all the ingredients required for crafting a pogostick appears in the grid world, and a new action ``approach plastic chest'' is added to the agent's action set. The chest can be acted upon with the ``collect'' action, upon which the ingredients contained in the chest are added to the agent's inventory.  \tabularnewline
    \arrayrulecolor{taupegray} \midrule
    
    \emph{Distance} & \emph{Detrimental} & One cell must separate the trader and the agent in order for the trade action to execute successfully. \tabularnewline
    \arrayrulecolor{taupegray} \midrule
    
    \emph{Fence} & \emph{Detrimental} & A breakable fence surrounds each tree in the grid world, and a new action ``approach fence'' is added to the agent's action set. \tabularnewline
    \arrayrulecolor{taupegray} \midrule
    
    \emph{Fire} & \emph{Detrimental} & The crafting table in the grid world is given a new property ``on fire'', and this property is set to true at the beginning of the episode. Somewhere else in the grid world, a water bucket is placed. The agent can extinguish the fire by choosing either the ``use'' action or the ``collect'' action while near the crafting table and holding the water bucket, whereafter it is returned a ``bucket'' in its inventory. The crafting table can only be crafted at after the fire has been extinguished. \tabularnewline
    \arrayrulecolor{taupegray} \midrule
    
    \emph{Moving} & \emph{Nuisance} & The traders are given action sets with ``move forward'', ``move backward'', ``move left'', and ``move right'' and select one of these actions randomly at every time step. \tabularnewline
    \arrayrulecolor{taupegray} \midrule
    
    \emph{Multi-interact} & \emph{Detrimental} & There are two variants of this novelty. Either trees cannot be broken with the ``break'' action or there are no trees in the grid world at all. In either case, the agent can obtain oak logs by interacting with the traders. \tabularnewline
    \arrayrulecolor{taupegray} \midrule
    
    \emph{Multi-room} & \emph{Nuisance} & Two rooms separated by a wall and doorway comprise the grid world rather than just one large room. \tabularnewline
    \arrayrulecolor{taupegray} \midrule
    
    \emph{Portal Treasure} & \emph{Detrimental} & The agent is given a treasure in its inventory, a portal is placed in the grid world, and no tree can be broken using the break action. If the agent selects the action ``use'' while at the portal and with the treasure in its inventory, it loses the treasure but receives a number of oak logs in its inventory instead. \tabularnewline
    \arrayrulecolor{taupegray} \midrule
    
    \emph{Random drop} & \emph{Detrimental} & The agent can no longer break trees using the break action nor any other action. However, with a certain probability, the action ``break block'' results in a given number of oak logs being dropped into the grid world. \tabularnewline
    \arrayrulecolor{taupegray} \midrule
    
    \emph{Space Around} & \emph{Detrimental} & At the beginning of the episode, the agent is given a number of saplings in its inventory. These saplings can be placed in the grid world and grow into trees. The agent cannot place a sapling when there is a tree or wall within radius 1 of the item. \tabularnewline
    \arrayrulecolor{taupegray} \bottomrule
    
    \end{tabular}
    \caption{Novelties currently implemented in the \NovelGym{} environment}
    \label{tab:rl_systemic_examples}
\end{table}


\subsection{Reward Function}
\label{sec:ng_reward_function}

\rev{The reward function plays a critical role in guiding the RL agent toward completing the multi-step crafting task. Because the pogostick task requires the agent to achieve a sequence of subgoals in order---collecting resources, crafting intermediaries, trading, and finally crafting the pogostick---a sparse reward signal (rewarding only final task completion) leads to prohibitively slow learning. Instead, we use a \emph{plan-based shaped reward} that leverages the symbolic planner to decompose the task into an ordered queue of subgoals and assigns intermediate rewards as subgoals are achieved.}

\rev{The reward-shaping mechanism works as follows. The planner generates a plan $P$ for the current task, from which an ordered queue $\mathcal{Q}$ of subgoals is extracted. The subgoals used in the pogostick task are:}
\begin{enumerate}[leftmargin=*]
    \item \rev{\texttt{break \textlangle block of platinum\textrangle}}
    \item \rev{\texttt{trade block of titanium}}
    \item \rev{\texttt{break block of diamond}}
    \item \rev{\texttt{break tree}}
    \item \rev{\texttt{craft planks}}
    \item \rev{\texttt{craft stick}}
    \item \rev{\texttt{craft pogostick}}
\end{enumerate}

\rev{A reward is only given for the first unachieved subgoal in the queue. Once a subgoal is achieved for the first time, it becomes a \emph{past goal} and the next subgoal in the queue becomes the active target. Two additional mechanisms prevent reward exploitation:}
\begin{itemize}[leftmargin=*]
    \item \rev{\textbf{Reward decay}: If the agent re-accomplishes the immediately preceding past goal (e.g., by collecting an already-collected resource type), the reward for that subgoal is halved. This discourages the agent from repeatedly achieving easy subgoals for reward.}
    \item \rev{\textbf{Inventory check}: When a subgoal is achieved, the algorithm scans subsequent subgoals in the queue and checks them against the agent's current inventory. If any subsequent subgoal is already satisfied (e.g., the agent already has the required items from an earlier action), that subgoal is marked complete with zero reward and skipped.}
\end{itemize}

\rev{Algorithm~\ref{alg:reward_func} formalizes this training routine.}

\begin{algorithm}
\caption{Training Routine with Crafted Reward Function}
\label{alg:reward_func}
\begin{algorithmic}[1]
\STATE Plan $P$
\STATE Ordered set (queue) $\mathcal{Q}$ of subgoals derived from $P$
\STATE Set of past goals $\mathcal{G}_{\text{past}} \gets \emptyset$

\WHILE{$\mathcal{Q}$ is not empty}
    \STATE Let $g_{\text{current}}$ be the first element of $\mathcal{Q}$
    
    \IF{$g_{\text{current}}$ is achieved}
        \STATE $\mathcal{G}_{\text{past}} \gets \mathcal{G}_{\text{past}} \cup \{ g_{\text{current}} \}$
        \STATE Remove $g_{\text{current}}$ from $\mathcal{Q}$
        
        \IF{$g_{\text{current}}$ was the immediate predecessor in $\mathcal{G}_{\text{past}}$}
            \STATE $\text{Reward}(g_{\text{current}}) \gets \frac{1}{2} \times \text{Reward}(g_{\text{current}})$
        \ENDIF
        
        \FOR{each $g_{\text{subsequent}}$ in $\mathcal{Q}$}
            \IF{$g_{\text{subsequent}}$ is achieved}
                \STATE Remove $g_{\text{subsequent}}$ from $\mathcal{Q}$
                \STATE $\text{Reward}(g_{\text{subsequent}}) \gets 0$
            \ELSE
                \STATE \textbf{break}
            \ENDIF
        \ENDFOR
    \ELSE
        \STATE Update model with $\text{Reward}(g_{\text{current}})$
    \ENDIF
\ENDWHILE
\end{algorithmic}
\end{algorithm}

\rev{
\paragraph{Complexity.}
We analyze the overhead of the reward-shaping routine itself, excluding the cost
of the underlying policy optimization. Let $n = |\mathcal{Q}|$ denote the number
of subgoals derived from the plan (bounded by the plan length, $n \leq |P|$).
Assume the queue $\mathcal{Q}$ and the past-goal set $\mathcal{G}_{\text{past}}$
are represented so that achievement checks, membership tests, insertions, and
single-subgoal removals are $O(1)$ on average, and the immediate-predecessor
test is $O(1)$ via a stored pointer.

Each of the $n$ subgoals is removed from $\mathcal{Q}$ exactly once over the
entire run---either as $g_{\text{current}}$ when it is achieved or inside the
skip-ahead inner loop---and the reward-decay and inventory-skip steps are $O(1)$
per subgoal. The nested loop therefore does not make the bookkeeping quadratic:
a subgoal removed in the inner loop is never revisited.

\begin{itemize}
  \item \textbf{Time:} $O(n) = O(|P|)$ for the reward-shaping bookkeeping. The
  overall wall-clock cost is dominated by the underlying policy optimization,
  which the routine only augments with a constant amount of work per step.
  \item \textbf{Space:} $O(n) = O(|P|)$ for the subgoal queue $\mathcal{Q}$ and
  the past-goal set $\mathcal{G}_{\text{past}}$, each holding at most $n$
  subgoals (excluding the memory of the underlying RL model).
\end{itemize}
}


\subsubsection{Pre-Novelty Base RL}

\rev{During pre-novelty training, the agent receives the plan-based shaped reward described above. The planner generates a complete plan for the pogostick task, and the subgoal queue guides the agent through the correct sequence of resource collection and crafting steps.}

\subsubsection{Post-Novelty Transfer RL}

\rev{After novelty injection, the transfer RL agent uses a simplified reward: $+1000$ for achieving the goal of crafting a pogostick, and $-1$ for each time step. The plan-based shaping is not used because the planner's model may no longer be valid in the transformed environment.}

\subsubsection{Post-Novelty Hybrid Planning and Learning}

\rev{The hybrid agent's post-novelty reward combines goal completion ($+1000$ for crafting a pogostick), a per-step penalty ($-1$ per time step), and a planning-failure penalty ($-250$ if the agent reaches a state from which the planner cannot generate a valid plan). The planning-failure penalty discourages the RL component from exploring states that are irrecoverable from the planner's perspective.}


\subsection{Sample Configuration File}
\label{sec:ng_sample_config}

\rev{Here we provide a sample configuration file, which generates trees in the world, allows the user to move around in the world, break the trees, and craft planks from the trees:}
\begin{lstlisting}
---
actions:
  break_block:
    module: gym_novel_gridworlds2.contrib.polycraft.actions.Break
    step_cost: 3600
  move_forward:
    module: gym_novel_gridworlds2.contrib.polycraft.actions.SmoothMove
    direction: W
  move_backward:
    module: gym_novel_gridworlds2.contrib.polycraft.actions.SmoothMove
    direction: X
  move_left:
    module: gym_novel_gridworlds2.contrib.polycraft.actions.SmoothMove
    direction: A
  move_right:
    module: gym_novel_gridworlds2.contrib.polycraft.actions.SmoothMove
    direction: D
  craft:
    module: gym_novel_gridworlds2.contrib.polycraft.actions.Craft
  trade:
    module: gym_novel_gridworlds2.contrib.polycraft.actions.Trade
action_sets:
  main:
  - break_block
  - craft_planks
  - move_*
object_types:
  default: gym_novel_gridworlds2.contrib.polycraft.objects.PolycraftObject
  oak_log: gym_novel_gridworlds2.contrib.polycraft.objects.easy_oak_log.OakLog
map_size: [16, 16]
rooms:
  '2':
    start: [0, 0]
    end: [15, 15]
objects:
  oak_log:
    quantity: 5
    room: 2
    chunked: 'False'
entities:
  main_1:
    agent: gym_novel_gridworlds2.agents.KeyboardAgent
    name: entity.polycraft.Player.name
    type: agent
    entity: gym_novel_gridworlds2.contrib.polycraft.objects.PolycraftEntity
    action_set: main
    inventory:
      iron_pickaxe: 1
      tree_tap: 1
    id: 0
    room: 2
    max_step_cost: 100000
recipes:
  planks:
    input:
    - oak_log
    - '0'
    - '0'
    - '0'
    output:
      planks: 4
    step_cost: 1200
trades: {}
auto_pickup_agents:
- 0
\end{lstlisting}


\subsection{PDDL Specification}
\label{sec:ng_pddl}

\subsubsection{Pre-Novelty Domain}

\rev{The pre-novelty domain is defined as follows:}
\begin{lstlisting}[
  style=pddlstyle,
  caption={Pre-novelty PDDL domain for the Polycraft task.},
  label={lst:ng-pre-domain}
]
(define (domain polycraft_generated)

(:requirements :typing :strips :fluents :negative-preconditions :equality)

(:types 
    pickaxe_breakable - breakable
    hand_breakable - pickaxe_breakable
    breakable - placeable
    placeable - physobj
    physobj - physical
    actor - physobj
    trader - actor
    pogoist - actor
    agent - actor
    oak_log - log
    distance - var
    agent - placeable
    trader - placeable
    pogoist - placeable
    bedrock - placeable
    door - placeable
    safe - placeable
    plastic_chest - placeable
    tree_tap - placeable
    oak_log - hand_breakable
    diamond_ore - pickaxe_breakable
    iron_pickaxe - physobj
    crafting_table - placeable
    block_of_platinum - pickaxe_breakable
    block_of_titanium - placeable
    sapling - placeable
    planks - physobj
    stick - physobj
    diamond - physobj
    block_of_diamond - physobj
    rubber - physobj
    pogo_stick - physobj
    blue_key - physobj
)

(:constants 
    air - physobj
    one - distance
    two - distance
    rubber - physobj
    blue_key - physobj
)

(:predicates ;todo: define predicates here
    (holding ?v0 - physobj)
    (floating ?v0 - physobj)
    (facing_obj ?v0 - physobj ?d - distance)
    (next_to ?v0 - physobj ?v1 - physobj)
)


(:functions ;todo: define numeric functions here
    (world ?v0 - physobj)
    (inventory ?v0 - physobj)
    (container ?v0 - physobj ?v1 - physobj)
)

; define actions here
(:action approach
    :parameters    (?physobj01 - physobj ?physobj02 - physobj )
    :precondition  (and
        (>= ( world ?physobj02) 1)
        (facing_obj ?physobj01 one)
    )
    :effect  (and
        (facing_obj ?physobj02 one)
        (not (facing_obj ?physobj01 one))
    )
)

(:action approach_actor
    :parameters    (?physobj01 - physobj ?physobj02 - actor )
    :precondition  (and
        (facing_obj ?physobj01 one)
    )
    :effect  (and
        (facing_obj ?physobj02 one)
        (not (facing_obj ?physobj01 one))
    )
)

(:action break
    :parameters    (?physobj - hand_breakable)
    :precondition  (and
        (facing_obj ?physobj one)
        (not (floating ?physobj))
    )
    :effect  (and
        (facing_obj air one)
        (not (facing_obj ?physobj one))
        (increase ( inventory ?physobj) 1)
        (increase ( world air) 1)
        (decrease ( world ?physobj) 1)
    )
)


(:action break_holding_iron_pickaxe
    :parameters    (?physobj - pickaxe_breakable ?iron_pickaxe - iron_pickaxe)
    :precondition  (and
        (facing_obj ?physobj one)
        (not (floating ?physobj))
        (holding ?iron_pickaxe)
    )
    :effect  (and
        (facing_obj air one)
        (not (facing_obj ?physobj one))
        (increase ( inventory ?physobj) 1)
        (increase ( world air) 1)
        (decrease ( world ?physobj) 1)
    )
)

(:action break_diamond_ore
    :parameters    (?iron_pickaxe - iron_pickaxe)
    :precondition  (and
        (facing_obj diamond_ore one)
        (not (floating diamond_ore))
        (holding ?iron_pickaxe)
    )
    :effect  (and
        (facing_obj air one)
        (not (facing_obj diamond_ore one))
        (increase ( inventory diamond) 9)
        (increase ( world air) 1)
        (decrease ( world diamond_ore) 1)
    )
)

(:action select
    :parameters    (?prev_holding - physobj ?obj_to_select - physobj)
    :precondition  (and
        (>= ( inventory ?obj_to_select) 1)
        (holding ?prev_holding)
        (not (= ?obj_to_select air))
    )
    :effect  (and
        (holding ?obj_to_select)
        (not (holding ?prev_holding))
    )
)

(:action deselect_item
    :parameters    (?physobj01 - physobj)
    :precondition  (and
        (holding ?physobj01)
        (not (holding air))
    )
    :effect  (and
        (not (holding ?physobj01))
        (holding air)
    )
)

(:action place_sapling
    :parameters (?sapling - sapling ?log - log)
    :precondition (and
        (facing_obj air one)
        (holding ?sapling)
    )
    :effect (and
        (facing_obj ?log one)
        (not (facing_obj air one))
        (increase ( world ?log) 1)
        (decrease ( inventory ?log) 1)
    )
)

(:action place
    :parameters   (?physobj01 - placeable)
    :precondition (and
        (facing_obj air one)
        (holding ?physobj01)
    )
    :effect (and 
        (facing_obj ?physobj01 one)
        (not (facing_obj air one))
        (increase ( world ?physobj01) 1)
        (decrease ( inventory ?physobj01) 1)
    )
)


(:action collect_from_tree_tap
    :parameters (?actor - actor ?log - log)

    :precondition (and
        (holding tree_tap)
        (facing_obj ?log one)
    )
    :effect (and
        (increase ( inventory rubber) 1)
    )
)

; additional actions, including craft and trade
(:action craft_stick
    :parameters ()
    :precondition (and
        (>= ( inventory planks) 2)
    )
    :effect (and
        (decrease ( inventory planks) 2)
        (increase ( inventory stick) 4)
    )
)


(:action craft_planks
    :parameters ()
    :precondition (and
        (>= ( inventory oak_log) 1)
    )
    :effect (and
        (decrease ( inventory oak_log) 1)
        (increase ( inventory planks) 4)
    )
)


(:action craft_block_of_diamond
    :parameters ()
    :precondition (and
        (facing_obj crafting_table one)
        (>= ( inventory diamond) 9)
    )
    :effect (and
        (decrease ( inventory diamond) 9)
        (increase ( inventory block_of_diamond) 1)
    )
)


(:action craft_tree_tap
    :parameters ()
    :precondition (and
        (facing_obj crafting_table one)
        (>= ( inventory planks) 5)
        (>= ( inventory stick) 1)
    )
    :effect (and
        (decrease ( inventory planks) 5)
        (decrease ( inventory stick) 1)
        (increase ( inventory tree_tap) 1)
    )
)


(:action craft_pogo_stick
    :parameters ()
    :precondition (and
        (facing_obj crafting_table one)
        (>= ( inventory stick) 2)
        (>= ( inventory block_of_titanium) 1)
        (>= ( inventory diamond) 1)
        (>= ( inventory rubber) 1)
    )
    :effect (and
        (decrease ( inventory stick) 2)
        (decrease ( inventory block_of_titanium) 1)
        (decrease ( inventory diamond) 1)
        (decrease ( inventory rubber) 1)
        (increase ( inventory pogo_stick) 1)
    )
)


(:action trade_block_of_titanium_1
    :parameters ()
    :precondition (and
        (facing_obj entity_103 one)
        (>= ( inventory block_of_platinum) 1)
    )
    :effect (and
        (decrease ( inventory block_of_platinum) 1)
        (increase ( inventory block_of_titanium) 1)
    )
)

)


\end{lstlisting}


\subsubsection{Pre-Novelty Problem}

\rev{The pre-novelty problem is defined as follows:}
\begin{lstlisting}[
  style=pddlstyle,
  caption={Pre-novelty PDDL problem instance for the Polycraft task.},
  label={lst:ng-pre-problem}
]
(define (problem polycraft_problem)
(:domain polycraft_generated)
    (:objects 
        agent - agent
        trader - trader
        pogoist - pogoist
        bedrock - bedrock
        door - door
        safe - safe
        plastic_chest - plastic_chest
        tree_tap - tree_tap
        oak_log - oak_log
        diamond_ore - diamond_ore
        iron_pickaxe - iron_pickaxe
        crafting_table - crafting_table
        block_of_platinum - block_of_platinum
        block_of_titanium - block_of_titanium
        sapling - sapling
        planks - planks
        stick - stick
        diamond - diamond
        block_of_diamond - block_of_diamond
        rubber - rubber
        pogo_stick - pogo_stick
        blue_key - blue_key
        entity_0 - agent
        entity_103 - trader
        entity_102 - pogoist
    )

    (:init
        (= (world air) 178)
        (= (world bedrock) 60)
        (= (world crafting_table) 1)
        (= (world entity_102) 1)
        (= (world oak_log) 5)
        (= (world entity_0) 1)
        (= (world entity_103) 1)
        (= (world diamond_ore) 4)
        (= (world plastic_chest) 1)
        (= (world block_of_platinum) 4)
        (= (world agent) 0)
        (= (world trader) 0)
        (= (world pogoist) 0)
        (= (world door) 0)
        (= (world safe) 0)
        (= (world tree_tap) 0)
        (= (world block_of_titanium) 0)
        (= (world sapling) 0)
        (= (inventory iron_pickaxe) 1)
        (= (inventory tree_tap) 1)
        (= (inventory agent) 0)
        (= (inventory trader) 0)
        (= (inventory pogoist) 0)
        (= (inventory bedrock) 0)
        (= (inventory door) 0)
        (= (inventory safe) 0)
        (= (inventory plastic_chest) 0)
        (= (inventory oak_log) 0)
        (= (inventory diamond_ore) 0)
        (= (inventory crafting_table) 0)
        (= (inventory block_of_platinum) 0)
        (= (inventory block_of_titanium) 0)
        (= (inventory sapling) 0)
        (= (inventory planks) 0)
        (= (inventory stick) 0)
        (= (inventory diamond) 0)
        (= (inventory block_of_diamond) 0)
        (= (inventory rubber) 0)
        (= (inventory pogo_stick) 0)
        (= (inventory blue_key) 0)
        (facing_obj air one)
        (holding air)
    )

    (:goal (>= (inventory pogo_stick) 1))
)

\end{lstlisting}



\subsection{Hyperparameters}
\label{sec:ng_hyperparameters}

\subsubsection{PPO}

\rev{Table~\ref{tab:ppo-param} shows the hyperparameters for the PPO algorithm:}
\begin{table}[htbp]
    \centering
    \begin{tabular}{cc}
    \toprule
        Name &  Value  \\
    \midrule
        Hidden Dense NN Layers  & 256, 64\\
        Optimizer & Adam \\
        Learning Rate & $10^{-5}$ \\
        $\gamma$ (discount factor) & 0.99 \\
        $\epsilon$ (PPO clip) & 0.2 \\
    \bottomrule
    \end{tabular}
    \caption{PPO hyperparameters.}
    \label{tab:ppo-param}
\end{table}

\subsubsection{ICM}

\rev{Table~\ref{tab:icm-param} shows the hyperparameters for the Intrinsic Curiosity Module:}
\begin{table}[htbp]
    \centering
    \begin{tabular}{cc}
    \toprule
        Name &  Value  \\
    \midrule
        Feature Net Hidden Layers & 256, 64 \\
        Feature Dimension & 16 \\
        Learning Rate Scale & 1 \\
        Reward Scale & 0.01 \\
        Forward Loss Weight & 0.2 \\
        Learning Rate & $10^{-4}$ \\
    \bottomrule
    \end{tabular}
    \caption{ICM hyperparameters.}
    \label{tab:icm-param}
\end{table}

\subsubsection{Training Dynamics}

\rev{Table~\ref{tab:pre-n-train-hyper} shows the hyperparameters for the pre-novelty training routine:}
\begin{table}[htbp]
    \centering
    \begin{tabular}{cc}
    \toprule
        Name &  Value  \\
    \midrule
        Epoch Length     & 28800\\
        Train Frequency  & 2400 \\
        Maximum Episode Length & 1200 \\
        Parallel Env Thread Count & 8 \\
        Batch Size       & 128 \\
        Batches per Train & 20 \\
        Episodes per Test & 100 \\
    \bottomrule
    \end{tabular}
    \caption{Pre-novelty training routine hyperparameters.}
    \label{tab:pre-n-train-hyper}
\end{table}

\rev{Different epoch lengths and maximum episode length are used for the post-novelty situation:}
\begin{table}[htbp]
    \centering
    \begin{tabular}{cc}
    \toprule
        Name &  Value  \\
    \midrule
        Epoch Length     & 4800\\
        Train Frequency  & 800 \\
        Maximum Episode Length & 400 \\
        Parallel Env Thread Count & 4 \\
        Batch Size       & 128 \\
        Batches per Train & 8 \\
        Episodes per Test & 100 \\
    \bottomrule
    \end{tabular}
    \caption{Post-novelty training routine hyperparameters.}
    \label{tab:post-n-train-hyper}
\end{table}


\subsection{Novelty Implementation Example}
\label{sec:ng_novelty_example}

\rev{Here is an implementation of an example novelty, which makes the traders busy 50\% of the time (and they will not trade with the agent when they are busy).}

\subsubsection{Configuration File}

\rev{This is the content of the novelty configuration YAML file. It specifies that the novelty needs to be injected at episode 0, and updates the action module associated with the \texttt{trade} action to be the new \texttt{BusyTrade} module. It also supplies the parameter of the new trade module, the percentage of time when the traders are busy.}
\begin{lstlisting}
---
novelties:
  '0':
    actions:
      trade:
        module: novelties.evaluation1.busy_traders.trade_busy.BusyTrade
        busy_ratio: 0.5
\end{lstlisting}

\subsubsection{Modified Trade Action Module}

\rev{This is the content of the modified \texttt{Trade} action module, which inherits from the default module. The new module raises an error when the internally generated random number is below the 50\% threshold, and gives the error message saying that it is busy.}
\begin{lstlisting}[language=Python]
from gym_novel_gridworlds2.actions.action import PreconditionNotMetError
from gym_novel_gridworlds2.contrib.polycraft.actions import Trade

import numpy as np

class BusyTrade(Trade):
    def __init__(self, busy_ratio=0, *args, **kwargs):
        self.busy_ratio = busy_ratio
        super().__init__(*args, **kwargs)

    def do_action(self, agent_entity, target_type=None,
                  target_object=None, **kwargs):
        threshold = self.dynamics.rng.uniform(0, 1)
        if threshold < self.busy_ratio:
            raise PreconditionNotMetError(
                "Trader is busy. Please try again later.")
        return super().do_action(
            agent_entity, target_type, target_object, **kwargs)
\end{lstlisting}

\rev{More examples of novelty implementations can be found in the NovelGym repository\footnote{\url{https://github.com/tufts-ai-robotics-group/NovelGym}}.}
